\section{Realizing the Dirichlet argument in finite arithmetic}
\label{sec:dirichlet-arithmetic}

The preceding argument is intentionally written in familiar mathematical notation. This section specifies its arithmetic meaning and the finite constructions behind the analytic abbreviations. It is part of the human proof, not a claim that a Lean proof object has already been constructed.

\subsection{The final quantifier-free target}

Let $Y(q,B)$ be the integer in Theorem~\ref{thm:dirichlet-explicit}. For $n\geq1$ put
\[
 F(n)=n!,\qquad
 R(n;q,a)=\sum_{p=2}^{n}
 [\operatorname{prime}(p)] [p\bmod q=a\bmod q]\,(F(n)\mathbin{\mathrm{div}}p).
\]
Each contributing $p$ divides $F(n)$, so $R(n;q,a)/F(n)$ is exactly the finite reciprocal-prime sum. The literal arithmetic conclusion is
\begin{equation}
 q\geq1\land\gcd(a,q)=1
 \quad\Longrightarrow\quad
 R(Y(q,B);q,a)>(B+1)F(Y(q,B)).
 \label{eq:dirichlet-integer-target}
\end{equation}
This is a quantifier-free formula in primitive recursive symbols; take its universal closure in $q,a,B$.

Define $W(q,a,B)$ to be the least $p\leq Y(q,B)$ with
\[
 B<p,\quad \operatorname{prime}(p),\quad p\bmod q=a\bmod q,
\]
and return $0$ if the bounded search finds none. The same search is defined for invalid inputs, using a harmless default modulus $\max(q,1)$ when necessary. If no successful $p$ occurred, the sum in~\eqref{eq:dirichlet-integer-target} would involve at most $B$ terms, each at most one. It would be at most $BF(Y)$, a contradiction. The bounded-search invariant therefore proves the full positive witness statement for $W$. Neither $W$ nor $Y$ is executed at its theoretical bound by the accompanying diagnostics.

\subsection{Characters as bounded finite tables}

The reduced residues form a list of at most $q$ entries. Its multiplication, inverses and identity are found by searches bounded by $q$. Lagrange's theorem here is the finite coset partition: form the translates of a subgroup, delete duplicate lists by equality tests, and count disjoint equal-sized blocks. The subgroup generated by given entries is computed by at most $q$ successive closure steps.

With $h=\varphi(q)$, represent a character by an $h$-entry list of exponents in $\{0,\ldots,h-1\}$. Check all multiplication equations modulo $h$. The subgroup-extension argument in Section~\ref{sec:direct-dirichlet} proves that precisely $h$ of the $h^h$ candidate lists pass. Roots and extensions are selected by bounded search among those lists. Real-valuedness is the finite test that every exponent is $0$ or, when $h$ is even, $h/2$. Its negation is therefore also decidable.

Here is a direct construction of compatible complex names which avoids any prerequisite about irreducibility of cyclotomic polynomials. Define
\[
 A(x)=\int_0^x\frac{dt}{1+t^2}\quad(0\leq x\leq1),\qquad \pi=4A(1).
\]
Use rational rectangles for the integral. Its integrand has derivative of absolute value at most $2$, so the rectangle error is explicit. Integrating the finite geometric identity also gives the alternating series for $A(1)$, with remainder at most $1/(2m+1)$ after $m$ terms. The first eight terms give $\pi>3$, and the integral comparison gives $\pi<4$.

Set
\[
 R(x)=\exp(2iA(x))\frac{1-ix}{1+ix}.
\]
Its formal derivative vanishes: both logarithmic derivative terms cancel after multiplication by $1+x^2$. All denominators have modulus at least $1$. This calculation has a finite mesh proof. On $[0,1]$, the exponential factor has first derivative bounded by $2$ and second derivative bounded by $8$; the rational factor has corresponding bounds $2$ and $4$ and modulus $1$. Thus $|R''|\leq20$. Taylor's first-order cell error, obtained by two rectangle summations, is at most $10/N^2$ on each of $N$ equal cells. Telescoping gives $|R(1)-R(0)|\leq10/N$ plus the explicitly budgeted sample errors. At arbitrary precision this proves $R(1)=R(0)=1$. Consequently $\exp(2iA(1))=i$, and the exponential addition law gives $\exp(2\pi i)=1$.

Define $\zeta_h=\exp(2\pi i/h)$ for $h\geq1$. The finite series addition law and conjugation give $\zeta_h^h=1$ and $|\zeta_h|=1$. For $h\geq2$ and $1\leq j<h$, replace $j$ by $j_0=\min(j,h-j)$ and put $t=\pi j_0/h$. Then $0<t<2$. Pairing successive terms in the sine series gives
\[
 \sin t\geq t-\frac{t^3}{6}\geq\frac t3,
 \qquad
 |1-\zeta_h^j|=2\sin t\geq\frac{2\pi}{3h}>\frac2h.
\]
The paired terms are nonnegative because their successive absolute ratios are at most $4/((2m+2)(2m+3))<1$. The same inequality follows at any finite precision from the explicit series remainder. It proves both distinctness of the named roots and a uniform positive separation margin.

No equality test for arbitrary approximated complex numbers is used. Equality of character values and the real/nonreal split are decided by exponents modulo $h$. Multiplication of these tables agrees with multiplication of the constructed roots by the finite exponential identity. Cancellation in orthogonality now uses the displayed lower margin $2/h$, so it also has a direct rational error bound. This supplies the exact finite algebra and the compatible numerical names used in the Dirichlet argument.

\subsection{The specific analytic tails}

Here are cutoffs for the only unbounded sums and integrals in the proof. The precision parameter is $r\in\N$. Increase it by the explicit finite amplification budget of Section~\ref{sec:foundations} before combining quantities.

\begin{enumerate}
\item \textbf{The endpoint $L(1,\chi)$.} Partial summation gives
\[
 \left|L(1,\chi)-\sum_{n=1}^{N}\frac{\chi(n)}n\right|\leq\frac{2q}{N}.
\]
Take $N=2^{r+3}(q+1)$. This works uniformly for nonprincipal characters, using their bounded period sums.

\item \textbf{Integral coefficients.} Put
\[
 I_j(N)=\frac1{j!}\int_1^N A_\chi(t)t^{-3}(\log t)^j\,dt.
\]
From $(\log t)^j/j!\leq2^j\sqrt t$ we obtain
\[
 |I_j(\infty)-I_j(N)|\leq\frac{2q\,2^j}{3N^{3/2}}.
\]
For all $j\leq K$, the coarser common bound $4q2^K/N$ suffices for the tails of the $l_j$. The $u_j$ use the same calculation with $q$ replaced by $1$ and $A_\chi(t)$ replaced by $\{t\}$. Thus a common cutoff $N\geq2^{r+5}(q+1)2^K$ is more than sufficient. Split the integrals at every integer. On each resulting unit interval the partial-sum function is constant and the fractional-part function is affine.

\item \textbf{Integral quadrature.} For the integrands just listed, all values and first derivatives on $[1,N]$, away from integer break points, are bounded by the primitive recursive integer
\[
 J=2^{K+5}(q+1)(N+1)^{K+5}(K+5)!.
\]
This follows by differentiating $t^{-3}(\log t)^j$ and using $t^{-1}\leq1$, $\log t\leq N$, $j\leq K$; division by factorials only improves the bound. Use at least $2^{r+5}NJ$ equal subintervals on each unit interval, and evaluate at midpoints. The rectangle error is then less than $2^{-r-4}$. Sample-evaluation errors are summed over the total interval length and are budgeted separately. The identity
$\int_1^\infty t^{-3}(\log t)^jdt=j!/2^{j+1}$ follows by $j$ integrations by parts with the same explicit boundary-tail bounds; its finite recurrence is checked by induction on the bounded index $j$.

\item \textbf{The positive coefficients.} Equation~\eqref{eq:c-coefficient-tail} gives the tail $8\,4^jN^{-1/4}$. After multiplication by $(7/4)^j$ and summation through $K$, the total is bounded by
\[
 2\,7^{K+1}N^{-1/4}.
\]
The integer cutoff
\[
 N=2^{4(r+4)}7^{4(K+1)}
\]
makes it less than $2^{-r-3}$. Each retained term is a finite divisor sum and an explicitly approximated logarithmic power. Keep the positive square terms separately; they supply the lower estimate even before the rest of the series is evaluated.

\item \textbf{Euler logarithms.} At $s\geq1+2^{-T}$, discard primes above
\[
 N=2^{(r+T+4)2^T}
\]
and prime-power exponents above $r+6$. The first tail is at most
$2N^{-2^{-T}}/2^{-T}\leq2^{-r-3}$; the second is at most $3\cdot2^{-(r+6)}$. These estimates follow by summing positive geometric series and the decreasing-integrand tail. All exponents at the chosen integer cutoff can be bounded without solving a real comparison problem.

\item \textbf{Integrating the logarithmic derivative.} The proof supplies the uniform margin $m_0=\delta/(8q)$ on $[1,2]$, where $\delta=2^{-D}$. Differentiating the partial-summation integral once more gives $|L''(s)|\leq4q$ there. Hence
\[
 \left|\left(\frac{L'}L\right)'\right|
 \leq \frac{4q}{m_0}+\frac{4q^2}{m_0^2}
 \leq288q^4\delta^{-2}.
\]
For precision $2^{-r}$, a mesh with at least $2^{r+4}288q^4\delta^{-2}$ cells suffices for its quadrature error. Evaluate $L,L'$ by their partial-summation integrals, using the margin to budget division. Uniformly for $1\leq s\leq2$, the integral tail for $L'$ is at most
\[
 q\int_N^\infty t^{-s-1}(1+s\log t)\,dt
 =qN^{-s}\left(\log N+\frac2s\right)
 \leq\frac{q(\log N+2)}{N}\leq\frac{5q}{\sqrt N}.
\]
The choice $N=2^{2(r+6)}(q+1)^2$ makes this smaller than $2^{-r-3}$; the $L$ tail is smaller still. This proves the bound on the particular Euler logarithm, without testing a branch of the complex logarithm.
\end{enumerate}

For the geometric coefficient tails in $l_j$ and $h_j$, use their displayed finite partial sums and exact geometric identities. The following finite convolution argument supplies the coefficient identity $c_j=\varepsilon+h_j$ in detail. In particular, the polynomial evaluated at $7/4$ still has only $K+1$ terms: no divergent expansion is introduced by the arithmetic realization.

\subsection{A finite coefficient calculation, without differentiating a limit}
\label{sec:finite-jet-closure}

The identity $c_j=\varepsilon+h_j$ deserves an explicit arithmetic
justification.  The following calculation replaces coefficient comparison
and repeated differentiation of infinite series by finite convolution and
finite summation by parts.  Throughout this subsection logarithms stand for
the rational names already specified; each displayed identity is verified
to an arbitrarily assigned rational accuracy by a finite calculation.

Put
\[
 b_j(t)=\frac{(\log t)^j}{j!t^2},\qquad
 Z_j(N)=\sum_{n=1}^N b_j(n),\qquad
 M_j(N)=\sum_{n=1}^N\chi(n)b_j(n).
\]
The convention is $(\log1)^0=1$.  Let $Z_j$ and $M_j$ denote the corresponding
names with their stated tails.  The elementary exponential inequality gives
\[
 \frac{(\log t)^j}{j!}\leq4^j t^{1/4}\qquad(t\geq1),
\]
and hence, for both moment sums,
\begin{equation}
 |Z_j-Z_j(N)|,\ |M_j-M_j(N)|
       \leq\frac43\,4^jN^{-3/4}
       \leq2\,4^jN^{-3/4}.
 \label{eq:moment-series-tail}
\end{equation}
These are bounds on ordinary positive tails, so they apply to complex
characters without cancellation.  The coarser estimate
\[
 |Z_j|,\ |M_j|,\ |Z_j(N)|,\ |M_j(N)|\leq3\,2^j
\]
follows from $(\log t)^j/j!\leq2^j\sqrt t$ and
$\sum_{n\geq1}n^{-3/2}\leq3$.

\paragraph{Finite summation by parts identifies the integral coefficients.}
For each integer $N\geq1$ one has the finite identity
\begin{equation}
 M_j(N)=A_\chi(N)b_j(N)-\int_1^N A_\chi(t)b'_j(t)\,dt.
 \label{eq:finite-jet-partial-summation}
\end{equation}
Indeed, split at the integers and telescope
\[
 \sum_{n=1}^{N-1}A_\chi(n)\bigl(b_j(n+1)-b_j(n)\bigr).
\]
This is a bounded double-sum identity.  On each unit interval the formula
for $b'_j$ is checked by the elementary product rule and the finite
logarithm difference estimates.  Replacing its integral by rectangles
leaves the explicitly bounded quadrature error; thus this argument does
not presume an infinitary summation-by-parts theorem.

For $j\geq1$,
\[
 -b'_j(t)=t^{-3}\left(\frac{2(\log t)^j}{j!}
            -\frac{(\log t)^{j-1}}{(j-1)!}\right),
\]
and $-b'_0(t)=2t^{-3}$.  These are exactly the integrands defining $l_j$.
The endpoint term in~\eqref{eq:finite-jet-partial-summation} is bounded by
\[
 |A_\chi(N)b_j(N)|\leq q4^jN^{-7/4}.
\]
Together with the moment tail and the integral tails previously given,
this proves $M_j=l_j$ to any requested accuracy, with primitive recursive
cutoffs.

Apply the same finite identity with the constant sequence $1$ in place of
$\chi$; its partial-sum function is $\lfloor t\rfloor=t-\{t\}$.  One more
finite integration by parts gives
\begin{equation}
 Z_j(N)=\mathbf1_{j=0}+\int_1^N b_j(t)\,dt
                         +\int_1^N\{t\}b'_j(t)\,dt.
 \label{eq:finite-zeta-jet}
\end{equation}
Here the endpoint terms at $N$ cancel exactly.  The moment integral tends
to $1$ with an explicit bound:
\begin{equation}
 \int_1^N b_j(t)\,dt
     =1-\frac1N\sum_{i=0}^j\frac{(\log N)^i}{i!},
 \qquad
 0\leq\frac1N\sum_{i=0}^j\frac{(\log N)^i}{i!}
       \leq\frac{2^{j+1}}{\sqrt N}.
 \label{eq:finite-log-moment}
\end{equation}
The first identity is proved by the finite recurrence
$U_0(N)=1-1/N$ and
$U_j(N)=U_{j-1}(N)-b_j(N)N$ for $j\geq1$.
For comparison, the corresponding recurrence for
$V_j(N)=\frac1{j!}\int_1^N t^{-3}(\log t)^jdt$ is
\[
 V_0(N)=\frac{1-N^{-2}}2,\qquad
 V_j(N)=\frac12V_{j-1}(N)
             -\frac{(\log N)^j}{2j!N^2}.
\]
This proves the coefficient estimate based on
$V_j(\infty)=2^{-j-1}$ by a bounded-index recurrence, without induction on
a formula quantifying over all real precisions.

By the definition of $u_j$, equations~\eqref{eq:finite-zeta-jet} and
\eqref{eq:finite-log-moment} give
\begin{equation}
 Z_j=1+u_j.
 \label{eq:zeta-jet-identification}
\end{equation}
For $j=0$ the contribution $\mathbf1_{j=0}$ is the constant $1$ already
included in $u_0$; for $j>0$ it is absent.

\paragraph{A finite rectangle gives the product coefficients.}
For a fixed $j$ the binomial identity and
$\log(uv)=\log u+\log v$ give the finite convolution identity
\begin{align}
 \sum_{i=0}^jZ_i(N)M_{j-i}(N)
  &=\sum_{u=1}^N\sum_{v=1}^N
       \frac{\chi(v)}{u^2v^2}
       \frac{(\log(uv))^j}{j!}.
 \label{eq:finite-moment-convolution}
\end{align}
The terms with $uv\leq N$ are exactly
\[
 c_j(N)=\sum_{n=1}^N\frac{a_n(\log n)^j}{j!n^2},
 \qquad a_n=\sum_{v\mid n}\chi(v).
\]
Every factor pair of $n\leq N$ occurs in the rectangle, so this last
assertion is an equality of bounded divisor sums.  The remaining terms have
$N<uv\leq N^2$, and their absolute sum is at most
\[
 \sum_{n>N}\frac{d(n)(\log n)^j}{j!n^2}
       \leq8\,4^jN^{-1/4}.
\]
The same bound controls $|c_j-c_j(N)|$.  Thus taking the explicitly bounded
tails in~\eqref{eq:finite-moment-convolution} proves
\begin{equation}
 c_j=\sum_{i=0}^j Z_i l_{j-i}
     =\sum_{i=0}^j u_i l_{j-i}+\sum_{i=0}^j l_i.
 \label{eq:finite-product-coefficient-identity}
\end{equation}
No rearrangement of an unbounded double series is needed: the exact
identity is on an $N$-by-$N$ rectangle, and all its discrepancies have
displayed positive tail bounds.

For clarity, one common bound is available.  To establish these identities
simultaneously for $j\leq K$ with final accuracy $2^{-r}$, first put
$p=r+2K+10$ and take
\[
 N=2^{4(p+K+8)}(q+1)^4 7^{4(K+1)}.
\]
Each of the unweighted moment tails, coefficient-integral tails, boundary
terms, and the bound in~\eqref{eq:finite-log-moment} is less than
$2^{-p-7}$.  The sum of the positive coefficient tails after weighting by
$(7/4)^j$ is also less than $2^{-p-7}$.  Passing from the truncated
moments to one convolution amplifies its individual tail errors by at most
$6(K+1)2^K\leq2^{2K+4}$.  These allowances are therefore smaller than
$2^{-r}$, with room for the additional rational sample errors.  Assign
those sample errors by the finite-operation budget before doing the
calculation.  Separate smaller cutoffs are possible, but unnecessary.

\paragraph{The endpoint coefficient sum.}
The remaining equality $\varepsilon=\sum_{i\geq0}l_i$ can also be checked
without a termwise interchange of unbounded operations.  Let
$E_H(x)=\sum_{i=0}^H x^i/i!$, with $E_{-1}=0$.  A finite sum of the integral
coefficient formulas gives
\[
 \sum_{i=0}^H l_i
 =\int_1^\infty A_\chi(t)t^{-3}
              \bigl(2E_H(\log t)-E_{H-1}(\log t)\bigr)dt.
\]
For $t\geq1$ both exponential polynomials are between $0$ and $t$.
Consequently the tail of this whole finite sum of integrals after $N$ is
at most $3q/N$, independently of $H$.  On $[1,N]$, choose $H$ large enough
that both $E_H$ and $E_{H-1}$ approximate $\exp$ to the assigned accuracy
on $[0,\log N]\subseteq[0,N]$.  The explicit exponential cutoff in the
foundations section supplies such an $H$, with $H$ at least any previously
specified coefficient index.  The integrand then approximates
$A_\chi(t)t^{-2}$, whose integral is $L(1,\chi)$ and whose tail is at most
$q/N$.  Both the interval and polynomial degree have thus been chosen by
primitive recursive bounds.

For a fixed $J$, split the finite sum at $J$ and use
$|l_i|\leq2q2^{-i}$ on its remaining coefficients.  The result at every
requested accuracy is
\[
 \left|\varepsilon-\sum_{i=0}^J l_i\right|
       \leq2q2^{-J}.
\]
It follows that $b_j=\sum_{i=0}^jl_i-\varepsilon$ has the bound
$|b_j|\leq2q2^{-j}$.  Substituting this definition into
\eqref{eq:finite-product-coefficient-identity} gives precisely
$c_j=\varepsilon+h_j$.

\paragraph{The induction invariant.}
Fix $q,K,r$ and the displayed cutoffs before beginning any of these loops.
The entries of the arrays are Gaussian rationals approximating
$Z_i(N),M_i(N),u_i,l_i$ and the partial integrals.  An invariant records
the stated finite sum or recurrence and its rational error allowance for
every index $i$ up to the current index.  This bounded conjunction is a
Boolean product.  The recurrence for the $V_i$ multiplies the previous
error by $1/2$ and adds the current local error; convolution has the finite
amplification bound just displayed.  The same finite-operation budget
handles powers of approximated logarithms.

Thus induction is used on equations and inequalities between coded
rationals, with bounded index assertions.  The accuracy $r$ is an external
number parameter in that induction, not an unbounded quantifier inside
its formula.  The resulting universally quantified error bounds are
obtained only after that quantifier-free induction is complete.  At the
final Dirichlet inequality the reserved positive slack lets one choose a
single sufficiently fine rational accuracy; there is no need to decide
equality of an arbitrary real name.


\subsection{Why quantifier-free induction suffices for these steps}

Fix the inputs $q,a,B$ and the accuracy assigned to each occurrence. Every character table, factor list, convolution, quadrature mesh, series truncation and error estimate now has a primitive recursive bound. The induction invariants are the partial-sum identities, finite product identities, and rational inequalities for all indices up to the current bound. Compile those bounded index assertions into Boolean products. Positivity is preserved by addition and multiplication of nonnegative rationals; the few divisions have the supplied margins.

The ordinary calculus notation can therefore be eliminated by the displayed finite identities and their error allowances. At the final stage the lower estimate has far more than a unit of slack. If necessary, use the rational-gap lemma for the rational endpoint~\eqref{eq:dirichlet-integer-target}: a putative failure gives an explicit rational gap and thus a finite precision contradicting the approximating estimates. This is an intuitionistic argument over decidable rational inequalities.

This describes a derivation in the intended first-order theory at the usual level of a human mathematical proof. Producing all the corresponding \texttt{PR} codes, their syntactic correctness lemmas, and the \texttt{Prf PRAAx} term is the formalization work laid out later. A claim that a compiler has already generated that term would be false.

\subsection{A supplementary proof-theoretic check}

There is also a useful independent way to check the qualitative conclusion. For a quantifier-free primitive recursive matrix $R$, a classical PRA proof of $\forall\mathbf x\exists y\,R(\mathbf x,y)$ yields a fixed primitive recursive witness term with an intuitionistic PRA verification. Here is the relevant mechanism.

Let $U$ be the universal closures of quantifier-free theorems of $\ipra$. Given a quantifier-free $A(n,\mathbf z)$, bounded search for its first false instance through $n$ proves in $\ipra$ the local least-counterexample property
\[
 A(0,\mathbf z)\land\neg A(n,\mathbf z)
 \ \longrightarrow\ b>0\land A(b-1,\mathbf z)\land\neg A(b,\mathbf z).
\]
Its universal closure belongs to $U$. Together with $A(0)$ and the induction-step hypothesis, it yields the induction axiom. Thus every fixed PRA proof can be viewed as a classical first-order proof from finitely many universal axioms in $U$. Herbrand's theorem gives finitely many witness terms whose quantifier-free instances have a provable disjunction. All their atomic tests are decidable in $\ipra$, so the finite propositional reasoning is available intuitionistically. Selecting the first successful term gives the required primitive recursive witness.

Applying this observation to the published PRA provability of Dirichlet infinitude gives a second metatheoretic route to the qualitative $\ipra$ conclusion. It does not give the explicit bound above, and no such transfer is asserted for the higher-quantifier density targets.
